\section{Sampling a dense completion block}\label{sec:dense}

The completion distributions required by the small-state chain have large
minimum margins, but their margins need not have comparable sizes.  We now
give the sampling subroutine, using the global parameters in
Equation~\eqref{eq:parameters}.  All random choices in its implementation
use finitely many independent unbiased bits.

\begin{theorem}[Dense completion draw]\label{thm:dense-draw}
Let a rectangular block have at most $m$ rows and $n$ columns, nonnegative
integer margins of equal total less than $C$, and every margin at least
$L=d^{12}$.  For each integer $h\ge1$, there is an algorithm that always
outputs a feasible integer table and whose law is within $2^{-h}$ in
total variation of the uniform law on the block's feasible tables.
Its worst-case bit running time is polynomial in $d,b,h$.  Empty blocks,
and blocks having a single row or a single column, have deterministic
completion algorithms with the same guarantee.
\end{theorem}

The proof first groups integer tables into scaled bins and gives a
constant lower bound on rejection-sampling acceptance. A log-concave
continuous density controls the bin normalizer and the conductance of the
finite bin chain. We then implement its transitions and the within-bin offsets
with bounded numbers of random bits.
The scaled-coordinate and softened-density construction follows the
geometric approach of Dyer, Kannan, and Mount~\cite[Sections~4--5]{DKM1997};
expanded-polytope rounding is also central to Morris~\cite[Sections~1.3--2]{Morris2002}.
The estimates and finite-bit sampler needed here are proved below.

\subsection{Scaled coordinates and rejection from bins}

The deterministic cases require no argument beyond the prescribed
margins, so suppose the block has $a,c\ge2$ rows and columns.  Write its
margins as $R_i,P_j$.  Reorder the rows and columns so that $R_a$ and $P_c$
are maximum, breaking ties by their original indices, and restore the
original order when returning a table.  The entries with $i<a,j<c$
determine the entire table.  We group their possible integer values into
intervals of lengths
\begin{equation}\label{eq:dense-scales}
 \begin{gathered}
 s_{ij}=\min(R_i,P_j)\quad(1\le i\le a,\ 1\le j\le c),\\
 B=d^8,\quad K_0=4B,\quad A=d^4,\\
 w_{ij}=\left\lfloor\frac{s_{ij}}B\right\rfloor
 \quad(i<a,\ j<c).
 \end{gathered}
\end{equation}
There are $e=(a-1)(c-1)\le d$ free entries.  Since $s_{ij}\ge L\ge2B$,
\begin{equation}\label{eq:width-bounds}
 \frac{s_{ij}}{2B}\le w_{ij}\le\frac{s_{ij}}B.
\end{equation}
Each nonnegative integer free entry has a unique representation
$X_{ij}=w_{ij}v_{ij}+\beta_{ij}$ with integer $v_{ij}\ge0$ and
$0\le\beta_{ij}<w_{ij}$.  For a feasible table, $X_{ij}\le s_{ij}$,
so $v_{ij}\le2B<K_0$.  Thus all feasible tables are represented by bins
indexed by
\[
 \mathcal V=\{0,\ldots,K_0-1\}^e
\]
and an offset in each free coordinate.  The number of bins depends only
on the dimensions, although an offset range can depend on the margins.
We will give weight one to every bin containing a feasible table and
choose the other weights by penalizing violations of nonnegativity.
An ideal proposal then samples a weighted bin and uniform offsets,
rejecting an infeasible table.

To define these weights, pass temporarily to scaled real coordinates.
For $u\in\R^e$, define the affine table $X(u)$ by
\begin{align*}
 X_{ij}(u)&=w_{ij}u_{ij} &&(i<a,\ j<c),\\
 X_{ic}(u)&=R_i-\sum_{j<c}w_{ij}u_{ij} &&(i<a),\\
 X_{aj}(u)&=P_j-\sum_{i<a}w_{ij}u_{ij} &&(j<c),\\
 X_{ac}(u)&=R_a-\sum_{j<c}P_j+
                    \sum_{i<a,\,j<c}w_{ij}u_{ij}.
\end{align*}
These formulas impose all the margins.  Let $T_{ij}=X_{ij}/s_{ij}$,
including the dependent entries, and set
\[
 \mathcal P=\{u\in\R^e:T_{ij}(u)\ge0\text{ for every }i,j\}.
\]
Every free scale contributing to a dependent entry is at most that
entry's scale.  For example, $s_{ij}\le s_{ic}$ follows from
$P_j\le P_c$, $s_{ij}\le s_{aj}$ from $R_i\le R_a$, and
$s_{ij}\le s_{ac}$ follows from both inequalities.  Thus every coefficient
in each normalized affine entry has absolute value at most $1/B$, and
there are at most $d$ of them.  Consequently
\begin{equation}\label{eq:dense-lipschitz}
 |T_{ij}(u)-T_{ij}(v)|\le\frac dB\norm{u-v}_\infty.
\end{equation}
Feasibility and Equation~\eqref{eq:width-bounds} imply
$\mathcal P\subset[0,2B]^e$.  The integer representation above belongs
to the closed unit bin $Q_v=v+[0,1]^e$, at the point with coordinates
$u_{ij}=v_{ij}+\beta_{ij}/w_{ij}$.

The feasible polytope has a point $u_*$ with uniform positive slack:
\begin{equation}\label{eq:dense-interior}
 T_{ij}(u_*)\ge\frac1{2d}.
\end{equation}
To construct it, first assign $s_{ij}/(2d)$ to every cell.  This uses at
most $cR_i/(2d)\le R_i/2$ in row $i$ and at most
$aP_j/(2d)\le P_j/2$ in column $j$.  The remaining nonnegative real
margins have equal totals and can be completed by the greedy
supply--demand construction.  The resulting table determines $u_*$.
The strict slack and Equation~\eqref{eq:dense-lipschitz} also show that
$\mathcal P$ has positive $e$-dimensional volume.

On the cube $[0,K_0]^e$ define
\begin{equation}\label{eq:dense-penalty}
 p(u)=\max\bigl(0,\max_{ij}(-T_{ij}(u)-d/B)\bigr),
 \qquad \rho(u)=2^{-Ap(u)}.
\end{equation}
The allowance $d/B$ makes $p$ vanish at the lower corner of every bin
meeting $\mathcal P$, by Equation~\eqref{eq:dense-lipschitz}.
The function $p$ is convex and $(d/B)$-Lipschitz in the sup norm;
therefore $\rho$ is positive, continuous, and log-concave.  Comparison
with $u_*\in[0,2B]^e$ gives $0\le p(u)\le4d$ on the cube.
Assign the dyadic bin weights
\begin{equation}\label{eq:dense-bin-weights}
 E_v=\lfloor Ap(v)\rfloor,\qquad W(v)=2^{-E_v},
 \qquad Z_B=\sum_{v\in\mathcal V}W(v).
\end{equation}
In particular $W(v)=1$ whenever $Q_v$ contains a feasible point, even
on its boundary.  For $u\in Q_v$,
\[
 \bigl|Ap(u)-\lfloor Ap(v)\rfloor\bigr|
 \le 1+Ad/B=1+d^{-3}<2,
\]
so
\begin{equation}\label{eq:bin-density-comparison}
 \tfrac14W(v)\le\rho(u)\le4W(v).
\end{equation}
This comparison lets integrals of $\rho$ control sums of bin weights;
overlapping bin faces have zero volume in these integrals.

An \emph{ideal trial} draws $v$ with probability $W(v)/Z_B$, then draws
independent uniform offsets $\beta_{ij}\in\{0,\ldots,w_{ij}-1\}$ and sets
\begin{equation}\label{eq:dense-offset-output}
 X_{ij}=w_{ij}v_{ij}+\beta_{ij}\qquad(i<a,\ j<c).
\end{equation}
It fills the remaining entries to give the prescribed margins and
accepts exactly when all entries are nonnegative.

\begin{proposition}[Bin rejection]\label{prop:bin-acceptance}
An ideal trial succeeds with probability at least $1/64$.  Conditional on
success, its output is exactly uniform on the feasible integer tables.
\end{proposition}
\begin{proof}
The unique bin-and-offset representation of a feasible integer table
has a bin of weight one.  The table's probability before rejection is
therefore
\begin{equation}\label{eq:dense-equal-table-mass}
 \frac1{Z_B}\prod_{i<a,\,j<c}\frac1{w_{ij}},
\end{equation}
independent of the table.  The dependent entries are integers because
they are integer sums and differences.  This proves conditional
uniformity, including tables on faces and widths of different sizes.

For the success bound, let $n_{\rm full}$ be the number of bins whose
entire closure is contained in $\mathcal P$.  Each has weight one and
accepts every offset, so its contribution to the success probability is
$1/Z_B$.  We compare $n_{\rm full}$ and $Z_B$ with $\vol(\mathcal P)$.
Set
$\mathcal P^- =\{u:T_{ij}(u)\ge2d/B\text{ for every }i,j\}$.
The affine map
$u\mapsto(1-\alpha)u+\alpha u_*$ with $\alpha=4d^2/B$
sends $\mathcal P$ into $\mathcal P^-$ by
Equation~\eqref{eq:dense-interior}.  Hence Bernoulli's inequality gives
\begin{equation}\label{eq:dense-inner-volume}
 \vol(\mathcal P^-)\ge(1-4d^2/B)^e\vol(\mathcal P)
 \ge(1-4d^3/B)\vol(\mathcal P)
 \ge\tfrac12\vol(\mathcal P).
\end{equation}
Except on grid boundaries, every point of $\mathcal P^-$ lies in a bin
whose entire closure is in $\mathcal P$: the change of any normalized
entry across that bin is at most $d/B$.  These bins cover $\mathcal P^-$
up to a null set and each has volume one, so
$n_{\rm full}\ge\vol(\mathcal P^-)\ge\vol(\mathcal P)/2$.

To bound $Z_B$, for any $\gamma\ge0$ consider the entire relaxed
polytope $\mathcal P_\gamma=\{u:T_{ij}(u)\ge-\gamma\text{ for all }i,j\}$
in $\R^e$, without a cube restriction.  For $u\in\mathcal P_\gamma$ the
point
\[
 z=\frac{u+2d\gamma u_*}{1+2d\gamma}
\]
belongs to $\mathcal P$, since the affine coefficients sum to one and
each $T_{ij}(z)\ge0$.  Thus
$\mathcal P_\gamma\subset
u_*+(1+2d\gamma)(\mathcal P-u_*)$, and
\begin{equation}\label{eq:dense-relaxed-volume}
 \vol(\mathcal P_\gamma)
 \le(1+2d\gamma)^e\vol(\mathcal P)
 \le\exp(2d^2\gamma)\vol(\mathcal P).
\end{equation}
This argument also covers points with negative free coordinates.
The sublevel set $\{p\le(j+1)/A\}$ inside the sampling cube is contained
in $\mathcal P_{d/B+(j+1)/A}$.  Summing the upper bounds for the penalty
layers and using Equation~\eqref{eq:bin-density-comparison} gives
\begin{align*}
 Z_B&\le4\int_{[0,K_0]^e}\rho(u)\,du\\
 &\le4\sum_{j\ge0}2^{-j}
       \exp\bigl(2d^2(d/B+(j+1)/A)\bigr)\vol(\mathcal P)
 \le32\vol(\mathcal P).
\end{align*}
For the last numerical bound put $a_0=2/d^2$ and $b_0=2/d^5$.
Since $a_0+b_0\le1/4$, both $\exp(a_0)$ and
$\exp(a_0+b_0)$ are at most $4/3$.  The geometric sum is at most
$4(4/3)/(1-(4/3)/2)=16$, which is stronger than asserted.
Consequently
\[
 \Pp\{\text{success}\}\ge\frac{n_{\rm full}}{Z_B}
 \ge\frac{\vol(\mathcal P)/2}{32\vol(\mathcal P)}=\frac1{64}.
\]
\end{proof}

\subsection{An integral interpolation inequality}

We supply the continuous inequality used to mix the bins.  It is an
analytic proof device; the algorithm never integrates or samples a real
density. This finite-box statement is a special case of the
Pr\'ekopa--Leindler inequality~\cite{Leindler1972,Prekopa1973}. We include
the one-dimensional transport argument and induction on slices to specify
the regularity needed here.

\begin{lemma}[Finite-box integral interpolation]\label{lem:integral-interpolation}
Let $S_1,S_2,S_0$ be bounded finite unions of open axis-aligned boxes in
$\R^q$, and let $0<t<1$, with
$(1-t)S_1+tS_2\subset S_0$.  Let $f_1,f_2,f_0$ be positive continuous
functions on a compact box containing the closures of these sets.
If
\[
 f_0((1-t)x+ty)\ge f_1(x)^{1-t}f_2(y)^t
 \quad(x\in S_1,\ y\in S_2),
\]
then
\begin{equation}\label{eq:integral-interpolation}
 \int_{S_0}f_0\ge
       \left(\int_{S_1}f_1\right)^{1-t}
       \left(\int_{S_2}f_2\right)^t.
\end{equation}
In particular, for finite unions $S,D_0$ of open boxes in the sampling
cube,
\begin{equation}\label{eq:dense-logconcave-interpolation}
 \int_{(1-t)S+tD_0}\rho\ge
       \left(\int_S\rho\right)^{1-t}
       \left(\int_{D_0}\rho\right)^t.
\end{equation}
\end{lemma}
\begin{proof}
We first prove the one-dimensional integral assertion in the slightly
larger class needed for slices.  Suppose $F,G,H$ are nonnegative,
supported on bounded finite unions of intervals, and, after removal of
finitely many endpoints, are continuous and either zero or bounded
above and below by positive constants on each interval piece.  Assume
$H((1-t)x+ty)\ge F(x)^{1-t}G(y)^t$ whenever $F(x)G(y)>0$, apart from
possible source endpoints.  Source integrals zero give a trivial
inequality; otherwise write $M_F=\int F>0$, $M_G=\int G>0$.

Let $u(a),v(a)$, $0<a<1$, be the quantiles of $F/M_F$ and $G/M_G$.
Split $(0,1)$ at the cumulative probabilities of the finitely many
support endpoints and continuity breakpoints of both densities.  On
each remaining open interval $I$, the quantiles are continuously
differentiable, strictly increasing, and satisfy
\[
 u'(a)=\frac{M_F}{F(u(a))},\qquad
 v'(a)=\frac{M_G}{G(v(a))}.
\]
These statements follow directly by inverting the continuously
differentiable cumulative function on a positive interval piece, whose
derivative is positive.  The function
$z(a)=(1-t)u(a)+tv(a)$ is also strictly increasing and continuously
differentiable on $I$.  Its images $z(I)$ for distinct pieces are
disjoint.  Change of variables on each piece, followed by the weighted
arithmetic--geometric mean inequality, yields
\begin{align*}
 \int H
 &\ge\sum_I\int_I H(z(a))z'(a)\,da\\
 &\ge\sum_I\int_I
 F(u(a))^{1-t}G(v(a))^t
 \left(\frac{M_F}{F(u(a))}\right)^{1-t}
 \left(\frac{M_G}{G(v(a))}\right)^t da\\
 &=M_F^{1-t}M_G^t.
\end{align*}
The pieces have total parameter length one.  A gap in a source support
causes a quantile jump at one of the excluded parameter values; the
corresponding gaps between image intervals merely omit nonnegative
contributions to $\int H$.  Thus neither disconnected supports nor
density jumps require continuity across a gap.

The assertion for dimension $q=1$ follows by restricting the ambient
continuous weights to their sets.  Induct on $q$.  For each first
coordinate $x$, let $S_i(x)$ be the slice in the remaining $q-1$
coordinates and put
\[
 F_i(x)=\int_{S_i(x)}f_i(x,x')\,dx'.
\]
Whenever $S_1(x)$ and $S_2(y)$ are nonempty, their interpolation is
contained in $S_0((1-t)x+ty)$.  The induction hypothesis for these slices
therefore gives
\[
 F_0((1-t)x+ty)\ge F_1(x)^{1-t}F_2(y)^t.
\]
To justify the one-dimensional step for the marginals, partition the
first-coordinate axis by all endpoints of all the boxes.  On each open
piece, a slice domain is a fixed union of boxes or is empty.  In the
nonempty case it has fixed positive volume.  Each ambient weight has a
positive minimum and finite maximum on the compact box, so its marginal
is bounded above and below by positive constants there.  Uniform
continuity of the weight on the compact box proves continuity of the
marginal on that piece by integration over this fixed domain.  These are
exactly the hypotheses of the one-dimensional argument just proved.
Applying that argument and then iterated integration proves
Equation~\eqref{eq:integral-interpolation}.

Finally, the interpolation of two finite unions of open axis-aligned
boxes is the finite union of all pairwise interpolations of their boxes,
each again such a box.  Convexity of $p$ gives the required pointwise
inequality for $f_0=f_1=f_2=\rho$, proving
Equation~\eqref{eq:dense-logconcave-interpolation}.
\end{proof}

\subsection{Conductance of the bin chain}

The following finite-chain estimate converts a cut bound into a mixing
bound. We will use it for the bin chain here and for the small-state
chain in Section~\ref{sec:algorithms}.

\begin{lemma}[Conductance and mixing]\label{lem:mixing}
Let $P$ be a reversible kernel on a nonempty finite set $V$, with
stationary probabilities $\pi(x)>0$ and $P(x,x)\ge1/2$.
Set $Q(x,y)=\pi(x)P(x,y)$ and
\[
 \mathcal E_P(H)=\sum_{\{x,y\}\subseteq V}
 Q(x,y)(H(x)-H(y))^2,
\]
where unordered pairs have distinct endpoints.  If
$\Var_\pi H\le K\mathcal E_P(H)$ for all $H$, with $K\ge1$, then,
writing $\pi_{\min}=\min_{x\in V}\pi(x)$, for every integer $t\ge0$,
\begin{equation}\label{eq:elementary-mixing}
 \norm{P^t(x,\cdot)-\pi}_{\TV}
 \le \pi_{\min}^{-1}e^{-t/K}.
\end{equation}
If, for some $\kappa>0$, every cut satisfies
\[
 Q(A,A^c):=\sum_{x\in A,\,y\notin A}Q(x,y)
 \ge\kappa\min\{\pi(A),\pi(A^c)\},
\]
then the variance inequality holds with $K=2/\kappa^2$.
\end{lemma}
\begin{proof}
For nonnegative $g$ with $\pi(g>0)\le1/2$, integrate the cut inequality
over $A_u=\{x:g(x)^2>u\}$ to obtain
\[
 \kappa\norm{g}_{L^2(\pi)}^2
 \le\sum_{\{x,y\}}Q(x,y)|g(x)^2-g(y)^2|
 \le\bigl(2\mathcal E_P(g)\norm{g}_{L^2(\pi)}^2\bigr)^{1/2}.
\]
The last step is Cauchy--Schwarz and
\[
 \sum_{\{x,y\}}Q(x,y)(g(x)+g(y))^2
 \le2\sum_x\pi(x)(1-P(x,x))g(x)^2
 \le2\norm{g}_{L^2(\pi)}^2.
\]
Consequently $\norm{g}_{L^2(\pi)}^2\le2\kappa^{-2}\mathcal E_P(g)$,
also when $g=0$.  Choose a $\pi$-median $a$ of $H$ and apply this to
$(H-a)_+$ and $(a-H)_+$.  Their squared norms sum to
$\norm{H-a}_{L^2(\pi)}^2$, and their energies sum to at most
$\mathcal E_P(H)$, by the corresponding inequality on each edge.
Since the variance is the minimum squared deviation from a constant,
this proves the cut assertion.

For mixing, reversibility makes $P$ self-adjoint on $L^2(\pi)$ and gives
$\mathcal E_P(H)=\langle H,(I-P)H\rangle_\pi$.
The stochastic kernel $2P-I$ is an $L^2(\pi)$ contraction by Jensen's
inequality and stationarity.  Hence the eigenvalues of $P$ are in $[0,1]$.
The variance inequality bounds those on the mean-zero subspace by
$1-1/K\le e^{-1/K}$.  For
$g_x(y)=\one_{\{y=x\}}/\pi(x)-1$, reversibility gives
\[
 (P^tg_x)(y)=\frac{P^t(x,y)}{\pi(y)}-1,
 \qquad \norm{g_x}_{L^2(\pi)}^2=\pi(x)^{-1}-1.
\]
The $L^2$ contraction followed by Cauchy--Schwarz bounds the total
variation distance by
\[
 \tfrac12\sqrt{\pi(x)^{-1}-1}\,e^{-t/K},
\]
which implies \eqref{eq:elementary-mixing}.
On a singleton all asserted variances and
distances vanish; it may simply be regarded as already mixed.
\end{proof}

Use the axis-neighbor Metropolis chain on $\mathcal V$.  Each of the
$2e$ signed coordinate directions is proposed with probability
\[
 q_B=2^{-q_{\rm bin}},\qquad
 q_{\rm bin}=\lceil\log_2(8d)\rceil;
\]
the unused probability, or a proposal leaving the cube, holds.  A
proposal $v'$ within the cube is accepted with probability
$\min(1,W(v')/W(v))$.  This chain is reversible for
$\pi_B(v)=W(v)/Z_B$ and holds with probability at least $3/4$.

\Needspace{8\baselineskip}
\begin{proposition}[Bin conductance and mixing]\label{prop:bin-gap}
For every nontrivial cut of the bin states, its stationary edge
conductance divided by the smaller stationary side mass is at least
\[
 \kappa_B=\frac1{10^5d^2K_0}.
\]
Starting at any bin, after
$H=d^{50}(h+1)^2$ steps the total-variation distance from $\pi_B$ is at
most $2^{-4(h+d)}$.
\end{proposition}
\begin{proof}
Write $\mu(E)=\int_E\rho$ and $D=(0,K_0)^e$.  For a cut
$\mathcal V=V_1\sqcup V_2$, choose $S$ to be the union of the open bins
on the side of smaller \emph{continuous} mass.  Then
$0<\mu(S)\le\mu(D)/2$.  Set
\[
 t=\frac1{4dK_0},\qquad r=tK_0=\frac1{4d},\qquad
 M=(1-t)S+tD.
\]
We have $S\subset M$, by writing $x=(1-t)x+tx$ for $x\in S$.
Lemma~\ref{lem:integral-interpolation} gives
\begin{equation}\label{eq:dense-mass-gain}
 \mu(M\setminus S)\ge(2^t-1)\mu(S)
                    \ge t\log(2)\mu(S).
\end{equation}
This is actual gained mass even though $S$ need not be convex.

We next cover this gain by neighborhoods of cut faces.  Ignore the
finite union of grid hyperplanes, a null set.  If $z\in M\setminus S$,
write $z=(1-t)x+ty$ with $x\in S$, $y\in D$.  The segment from $x$ to
$z$ stays inside the open convex cube and has sup-norm length less than
$r$.  Its initial bin lies on the selected side and its terminal bin
lies on the other side.  Follow the segment through its finitely many
grid crossings.  At a crossing in several coordinates simultaneously,
the incident bins are indexed by the two choices in each crossing
coordinate, and are connected by axis-neighbor faces through that same
point.  If the incoming and outgoing bins have different labels, a face
path between them contains an opposite-label pair.  Its common closed
face contains the crossing point.  Since some crossing changes the
label, $z$ is at sup distance less than $r$ from a cut face.  No exterior
face is needed: the entire segment lies in $D$.

For such a closed unit face $F$, take its sup-distance-$r$ neighborhood
intersected with $D$, denoted $F^{(r)}$.  Its volume is at most
\[
 2r(1+2r)^{e-1}
 \le\frac1{2d}\exp\left(\frac{e-1}{2d}\right)\le1.
\]
If $v$ is the lower corner of either bin adjacent to $F$, every point
$z\in F^{(r)}$ satisfies $\norm{z-v}_\infty\le1+r$.  Hence
\[
 \rho(z)\le2^{(Ad/B)(1+r)}\rho(v)\le4W(v),
\]
and the same inequality holds for the other adjacent corner $v'$.
Thus $\mu(F^{(r)})\le4\min(W(v),W(v'))$.  Combining the face cover with
Equation~\eqref{eq:dense-mass-gain}, and writing
$D_i=\sum_{v\in V_i}W(v)$, gives
\begin{align*}
 \sum_{\substack{vv'\text{ axis-neighbors}\\v\in V_1,\ v'\in V_2}}
             \min(W(v),W(v'))
 &\ge\frac{t\log(2)}4\mu(S)\\
 &\ge\frac{t\log(2)}{16}\min(D_1,D_2).
\end{align*}
The second inequality follows from
Equation~\eqref{eq:bin-density-comparison}: the selected continuous side
has mass at least one quarter of its own discrete weight, which is at
least $\min(D_1,D_2)$.  It does not require the two notions of smaller
side to agree.  Stationary edge conductance is
$q_B\min(W(v),W(v'))/Z_B$.  Since $q_B\ge1/(16d)$, the cut ratio is at
least $\log(2)/(1024d^2K_0)\ge\kappa_B$.

The weights satisfy $2^{-4Ad}\le W(v)\le1$, so
\[
 \pi_{B,\min}\ge\frac{2^{-4Ad}}{K_0^e},\qquad
 \log\pi_{B,\min}^{-1}\le4d^5\log2+d\log(4d^8)\le4d^5.
\]
Lemma~\ref{lem:mixing} supplies inverse gap at most
\[
 K_B=2\kappa_B^{-2}=3.2\cdot10^{11}d^{20}\le d^{32}
 \quad(d\ge14).
\]
It bounds the error at time $H$ by
$\exp(4d^5-H/K_B)$.  Since $d^{18}\ge4d^5$ and
$d^5h(h+2)\ge(d+1)h\ge h+d$, we have
\[
 \frac H{K_B}-4d^5\ge d^{18}(h+1)^2-4d^5
 \ge4d^5h(h+2)\ge4(h+d)\log2.
\]
This proves the asserted error.
\end{proof}

\subsection{A bounded random-bit implementation}

We finish the proof of Theorem~\ref{thm:dense-draw}, keeping explicit
bounds that will also describe the finite output law in
Section~\ref{sec:algorithms}.  Use the integer schedule
\begin{equation}\label{eq:dense-schedule}
 J=d^4(h+1),\qquad H=d^{50}(h+1)^2.
\end{equation}
Run at most $J$ independent trials, each running $H$ bin-chain steps from $0$.
At the end of each bin walk use Equation~\eqref{eq:dense-offset-output},
but obtain each offset from the bounded draw
\begin{equation}\label{eq:dense-offset-bits}
 \begin{gathered}
 k_{ij}=\lceil\log_2(w_{ij}+1)\rceil+4(h+d),\quad
 Z_{ij}\text{ uniform on }\{0,\ldots,2^{k_{ij}}-1\},\\
 \beta_{ij}=\left\lfloor\frac{w_{ij}Z_{ij}}{2^{k_{ij}}}\right\rfloor.
 \end{gathered}
\end{equation}
Return the first feasible table.  If all $J$ trials fail, return a fixed
greedy feasible table with the required margins.

Each offset probability differs from $1/w_{ij}$ by at most $2^{-k_{ij}}$:
its number of preimages is either the floor or the ceiling of
$2^{k_{ij}}/w_{ij}$.  Thus its total-variation error is at most
$w_{ij}/2^{k_{ij}}\le2^{-4(h+d)}$.  Coupling the independent offsets
coordinate by coordinate bounds their joint error by
$d2^{-4(h+d)}$.  Proposition~\ref{prop:bin-gap} bounds the bin error by
$2^{-4(h+d)}$, so a complete implemented trial, including its success
flag, differs from an ideal trial by at most
$(1+d)2^{-4(h+d)}$.

Couple the $J$ independent trial records until the first discrepancy.
This bounds the difference between the implemented first-success
procedure and its ideal counterpart by
$J(1+d)2^{-4(h+d)}$, without conditioning on an approximate success event.
By Proposition~\ref{prop:bin-acceptance}, the ideal procedure has uniform
law on success and failure probability at most $(63/64)^J$.  Its greedy
fallback therefore changes the uniform law by at most this probability.
The total error is at most
\begin{equation}\label{eq:dense-total-error}
 (63/64)^J+J(1+d)2^{-4(h+d)}\le2^{-h}.
\end{equation}
For explicit numerical bounds, $d^4\ge128$ gives
$(63/64)^J\le\exp(-J/64)\le2^{-2(h+1)}\le2^{-h-2}$.
Also $d\le2^{d/2}$, $d+1\le2^d$, and $h+1\le2^h$ give
$J(1+d)2^{-4(h+d)}\le2^{-3h-d}\le2^{-h-2}$.
All actual outputs are feasible by the acceptance test or the fallback.

It remains to account for every computation in the bit model.
At an integer bin corner $v$, the entries $X_{ij}(v)$ are integer sums
of products $w_{ij}v_{ij}$ and margins.  Since $w_{ij}<C/B$ and
$v_{ij}<4B$, their absolute values are at most a fixed multiple of $dC$;
they have $O(b+\log d)$ bits.  The candidates for the maximum defining
$p(v)$ are the rational numbers
\[
 0\quad\text{and}\quad
 \frac{-BX_{ij}(v)-ds_{ij}}{Bs_{ij}}.
\]
Cross multiplication, comparison, multiplication by $A$, and integer
division compute $E_v=\lfloor Ap(v)\rfloor$ exactly with polynomial bit
cost.  Denominators have $O(b+\log d)$ bits and there are at most $d$
candidates.  In particular no evaluation of an exponential or integral
is required.  We have $0\le E_v\le4Ad$, so a Metropolis acceptance is
either certain or has probability $2^{-(E_{v'}-E_v)}$, implemented by
requiring that many fresh unbiased bits to be zero.  The difference has
absolute value at most $8Ad$, a convenient uniform allowance.  A draw
of $q_{\rm bin}$ bits proposes each signed coordinate direction by one
code, with the remaining codes holding; this implements the stated
proposal exactly.

The offset draw in Equation~\eqref{eq:dense-offset-bits} uses
$O(b+h+d)$ bits and an integer multiplication and division on numbers
of polynomial bit length.  Table construction, feasibility tests, the
greedy fallback, and the bounded loop counters also use polynomially
many operations on polynomial-length integers.  There are $JH$
transitions and at most $dJ$ offsets.  Schoolbook integer arithmetic
already yields a worst-case polynomial bound in $d,b,h$.

For later use, a uniform upper bound on the number of bits per transition
and per offset is respectively
\begin{equation}\label{eq:dense-denominator-bounds}
 \begin{gathered}
 Q_0=q_{\rm bin}+8Ad,\qquad
 r_{\max}=b+1+4(h+d),\\
 R_1=HQ_0+dr_{\max},\qquad R_D=JR_1.
 \end{gathered}
\end{equation}
Unused transition bits, missing free coordinates, and unused trials
after the first success can be padded with unread independent bits.
Thus a trial has probabilities with common denominator $2^{R_1}$, and
the whole subroutine has probabilities with common denominator
$2^{R_D}$.  This padding changes no output law.  For an offset of width
$w$ and actual bit length $k$, its exact probability of value $\beta$ is
\begin{equation}\label{eq:dense-offset-count}
 2^{-k}\left(
 \left\lceil\frac{(\beta+1)2^k}{w}\right\rceil-
 \left\lceil\frac{\beta2^k}{w}\right\rceil\right),
 \qquad 0\le\beta<w.
\end{equation}
It can therefore be tabulated by integer arithmetic without enumerating
the random bit strings.  The deterministic completion cases can use the
same denominator allowances.  This completes the proof of
Theorem~\ref{thm:dense-draw}.
